%%%%%%%%%%%%%%%%%%%%%%%%%%%%%%%%%%%%%%%%%%%%%%%%%%%%%%%%%%%%%%%%%%%%%%
%  月度阅读总结 · 阶段一（初步阅读）  ——  陈思文 · 微通道液冷设计
%
%  编译：xelatex 文件名.tex   （跑两遍）
%     Overleaf：Menu -> Compiler -> XeLaTeX
%%%%%%%%%%%%%%%%%%%%%%%%%%%%%%%%%%%%%%%%%%%%%%%%%%%%%%%%%%%%%%%%%%%%%%
\documentclass[stage1]{ydzjreport}

% ===================== 只改这里：你的信息 =====================
\stageonesetup
  {阶段一 · 初步阅读}%          1 报告阶段
  {陈思文}%                   2 姓名
  {微通道液冷设计}%             3 方向（2-6 个汉字）
  {2026-07-01 至 2026-08-31}%   4 报告周期
  {微通道；仿生拓扑；Fluent操作；基础仿真}%   5 关键词
  {完成芯片散热前沿文献调研，掌握Fluent网格划分与基础工作流，并完成基础肋片与坑道换热的初步对比。}%          6 一句话摘要
% ============================================================

\begin{document}
\makestageonehead

% ==================================================================
\section{阅读范围}

\subsection{文献清单}

% 必填，至少 3 篇。列数固定 4 列，不要加列。
\begin{table}[htbp]
\centering
\caption{本月阅读文献}
\begin{tabularx}{\textwidth}{@{}c Y c Y@{}}
\toprule
\rowcolor{ydzjlight} 序号 & 文献 & 类型 & 与本方向的关系 \\
\midrule
1 & 袁永帅等 (2024)，《当前芯片散热技术方法和特点研讨》 & 期刊 & 明确芯片散热技术的宏观背景与工作温度限制 \\
2 & 熊园园等 (2024)，《基于内嵌式微通道芯片散热结构设计研究综述》 & 期刊 & 明确传统均匀通道存在“过度冷却”的痛点 \\
3 & 苗杰 (2024)，《面向服务器CPU芯片散热的变截面通道液冷板设计及性能分析》 & 学位论文 & 提供变截面通道设计的思路参考 \\
\bottomrule
\end{tabularx}
\end{table}

\subsection{选读理由}

阅读线索从宏观到具体：先通过综述类文献（文献1）了解当前芯片散热技术的整体格局、数据中心能耗指标要求以及芯片工作温度上限，明确“为什么要做高效液冷”；再通过内嵌式微通道综述（文献2）认识传统均匀微通道在间接冷却方式下的热流密度上限及其固有缺陷；最后聚焦到变截面通道液冷板的设计实例（文献3），从中寻找仿生微通道与变截面扰流结构的几何灵感，为后续翼型扰流柱的选型提供依据。


% ==================================================================
\section{核心概念梳理}

\subsection{概念一：局部热流密度突变与过度冷却}

芯片发热集中于结区，热流密度在空间上高度不均匀；若采用均匀布置的微针翅片或等截面直通道，无热源背景区会被“一视同仁”地强力冷却，造成过度冷却，形成过大的温度梯度并浪费泵功。对本项目而言，这意味着不能采用均匀结构，而必须按热流密度对通道进行“分区设计”，使换热能力与局部热负载相匹配。

\subsection{概念二：强迫对流与流体扰动}

强迫对流换热的关键在于破坏壁面热边界层。通过改变通道截面形状（如设置肋片、凹坑）可以诱发次级流动，使冷流体不断冲刷换热壁面，从而在相同泵功下获得更高的换热系数。对本项目而言，这意味着需要在微通道内布置翼型肋片或凹坑结构，利用其尾流扰动强化换热，同时兼顾由此带来的流阻代价。


% ==================================================================
\section{阶段一：通读与摘录}

\subsection{学习路径}

% 按时间或逻辑顺序说明「先看什么、再看什么、为什么」。3--6 条。
\begin{enumerate}
  \item 查找并研读芯片散热领域论文，学习前人在微通道散热方面的研究成果与结论。
  \item 从仿生结构（如月牙型沙丘、叶脉型分形）中寻找微通道扰流单元的结构灵感。
  \item 学习 Workbench Meshing 网格划分方法，掌握 Discovery 建模与 Fluent 求解器的搭配使用流程。
  \item 学习 UDF 文件导入，以及稳态/瞬态求解的参数设置与常见报错（负体积、不收敛等）的处理方法。
\end{enumerate}

\subsection{关键摘录}

% 必填，至少 3 条。第二列写原文要点，第三列写你自己的理解（不是复述）。
\begin{table}[htbp]
\centering
\caption{关键摘录与理解}
\begin{tabularx}{\textwidth}{@{}Y Y Y@{}}
\toprule
\rowcolor{ydzjlight} 文献 & 原文要点 & 我的理解 \\
\midrule
袁永帅等 (2024) & 数据中心能耗指标对芯片散热技术提出明确要求，芯片存在工作温度上限，散热能力直接制约性能释放 & 散热不是附属性问题，而是芯片性能与可靠性设计的一级约束，方案选型必须以温度上限为硬指标 \\
熊园园等 (2024) & 均匀内嵌式微针翅片在无热源背景区会造成过度冷却，产生过大的温度梯度 & 传统均匀方案存在“效能饱和”：一味加密翅片不仅收益递减，还浪费泵功，因此结构必须与热流密度分布耦合 \\
苗杰 (2024) & 变截面通道能够周期性破坏与重建热边界层，在适度压降代价下显著强化换热 & 不连续肋结构诱发的次级流动是打破效能饱和的关键，翼型截面兼具扰流与低阻力的潜力，值得作为基本扰流单元 \\
\bottomrule
\end{tabularx}
\end{table}


% ==================================================================
\section{阶段二：导出与复现}

% ★ 阶段一的核心：只读不做不算完成。
%   必须给出可复现的环境、步骤、结果。

\subsection{复现目标}

基础肋片结构与坑道（凹坑）结构的换热与压降效果对比测试：在相同入口条件与热流密度下，分别建立含肋片和含坑道的微通道模型，对比二者的温度峰值、壁面换热能力与进出口压降，验证“扰流强化换热但增大流阻”的基本规律。

\subsection{复现环境与步骤}

% 软件 + 版本 + 关键设置 + 步骤。要求：别人照着能重复出来。
\begin{enumerate}
  \item 环境：ANSYS Workbench（Discovery 建模 + Meshing 网格划分 + Fluent 求解），水为冷却工质，固体域与流体域耦合求解。
  \item 建模与网格：在 Discovery 中分别建立含肋片与含坑道的微通道几何，抽取流体域；在 Meshing 中设置网格尺寸控制与边界层加密，检查网格畸变度与正交质量，剔除负体积单元。
  \item 求解设置：采用稳态与瞬态求解器分别测试；设置入口边界条件（速度/质量流率）与壁面热流密度，监控残差、壁面温度与进出口压差直至收敛。
  \item 报错处理：针对求解过程中的负体积报错与收敛振荡，通过细化网格、调整松弛因子与时间步长予以排除。
  \item 结果导出：提取底板平均温度、最高温度、进出口压差及面积加权换热系数，整理成对比数据表。
\end{enumerate}


\subsection{复现结果}

% 必须给图或表。并说明与原文的差异在哪。
\begin{table}[htbp]
\centering
\caption{肋片与坑道结构仿真对比结果（定性/定量汇总）}
\begin{tabularx}{0.85\textwidth}{@{}Y c c c@{}}
\toprule
\rowcolor{ydzjlight} 项目 & 坑道结构 & 肋片结构 & 综合评价 \\
\midrule
流阻表现 & 较低，压降平缓 & 明显更高，随流速增长快 & 坑道占优 \\
换热表现 & 增强有限 & 壁面换热强化显著 & 肋片占优 \\
温度均匀性 & 一般 & 较好，温度峰值更低 & 肋片占优 \\
\bottomrule
\end{tabularx}
\end{table}

初步结论：肋片结构通过持续扰动热边界层获得了更强的换热能力与更低的温度峰值，但付出了更高的压降代价；坑道结构流阻低但换热增益有限。这与文献中“扰流强化换热伴随流阻上升”的结论一致，说明本次仿真流程与结果可信，为后续翼型扰流柱选型提供了方向。


% ==================================================================
\section{阶段性收获}

\subsection{认知上的变化}

我原以为均匀、密集布置的微通道散热效果最好，通道越密换热越强；现在认识到，均匀密集结构不仅容易在无热源区引发过度冷却、拉大芯片表面温度梯度，还会带来不必要的泵功损耗。散热结构设计必须针对热流密度分布进行“分区差异化”，在不同热流密度区域采用不同的扰流强度与结构形式。

\subsection{尚存的疑问}

% 至少 2 条，要具体到可以回答的问题。
\begin{enumerate}
  \item 对于中间热流密度变动极大的“浮动区”，应如何设计肋片构型，才能在热负载波动时既不欠冷也不过冷？
  \item 具体哪种翼型参数（相对厚度、相对弯度、攻角等）能够在换热强化与压降代价之间取得最优平衡？
\end{enumerate}


% ==================================================================
\section{下月计划}

\begin{table}[htbp]
\centering
\caption{下月工作计划}
\begin{tabularx}{\textwidth}{@{}Y Y c@{}}
\toprule
\rowcolor{ydzjlight} 事项 & 预期产出 & 时间 \\
\midrule
梳理并确立微通道液冷场景下的翼型参数 & 翼型参数候选方案 & 第 1 周 \\
开展不同翼型在低/高流速下的散热效率仿真 & 流热耦合仿真数据与对比分析 & 第 2--3 周 \\
开展六种不同构型的网格无关性检验 & 网格无关性检验报告与可靠基准数据 & 第 4 周 \\
\bottomrule
\end{tabularx}
\end{table}


% ==================================================================
% 参考文献：按正文引用顺序排列
% 正文引用写法：见表~\ref{tab:xxx}；文献引用用 \cite{key}
% 若不想用 \cite，也可直接写 [1]
\begin{thebibliography}{9}
\bibitem{ref1} 袁永帅, 等. 当前芯片散热技术方法和特点研讨[J]. 2024.
\bibitem{ref2} 熊园园, 等. 基于内嵌式微通道芯片散热结构设计研究综述[J]. 2024.
\bibitem{ref3} 苗杰. 面向服务器CPU芯片散热的变截面通道液冷板设计及性能分析[D]. 2024.
\end{thebibliography}

\end{document}
